\section{都城之间的地方，不是制度的空白}

938年幽州成为辽南京，1005年宋辽盟约将使节、边市与银绢编入关系，两个节点把城市、道路与边境财政放在同一张地图上。南京的城址、墓志、契丹与汉文题记可帮助识别某些官员、家族和宗教网络；它们不能证明每一座城或每一条商路都以相同方式被征收、供给或防卫。多都城并不是抽象制度，而是把不同资源区接到中枢的实践。

边市和盟约同样不能只用“朝贡”解释。宋辽双方的礼仪称谓、实际货物流动、堡寨维护与地方管控须分别检验；银绢的约定不等于边民的全部生计。辽的食货与地理志、宋人使辽记录、碑刻和考古各自照亮不同部分。承认它们的沉默，才能避免将燕云、上京与东北社会都写成同一笔财政账。

\subsection{南京与岁币的两端}

九三八年辽以幽州为南京，经营燕云城市与交通节点；一〇〇五年澶渊盟约将银绢、使节与边市编进宋辽关系。前者是把城市、道路和官署接入多都城体系，后者是以可重复输送的物资降低常态战争成本；两者都不等于边民生活被同一尺度的财政覆盖。银绢数额也不宜直接译为单向臣属关系。

《辽史》营卫、百官诸志、宋人使辽记录、盟约文本与墓志、城址能分别提示制度、礼仪与局部空间。它们只能让南京和边市的供给问题比本纪更清晰，不能据此量化草原、东北或燕云全部家庭的税役和交易。
